% \section{Discussion and Conclusion}
% FieldFormer offers a scalable and physically grounded framework for spatio-temporal field reconstruction, achieving strong performance across diverse benchmarks and surpassing both physics-regularized data-driven and classical baselines in accuracy. By combining local transformer inference with learnable anisotropic neighborhoods and differentiable physics regularization, the model enables mesh-free reconstruction on sparse, irregular sensor networks. More broadly, FieldFormer highlights how locality-aware inductive biases enable reliable knowledge extraction from sparse, longitudinal sensor data, a setting common in environmental monitoring, urban systems, and scientific sensing applications. While promising, some challenges remain: defining velocity-scaled local neighborhoods may be less effective in highly nonlocal or chaotic systems; autograd-based physics losses can be sensitive near sharp gradients or discontinuities; and the approach assumes that governing PDEs are at least partially known, whereas real-world dynamics may involve hidden processes. Moreover, predictive quality may diminish under extreme sparsity or sensor drift. Addressing these limitations through uncertainty-aware inference, robust physics integration, and hybrid discovery of unknown dynamics represents an important direction for future work.



% \subsection{Strengths, Limitations, and Diagnostic Insights}

% \noindent\textbf{Strengths: Locality-Aware Reconstruction and Parameter Efficiency.}
% FieldFormer is designed to recover coherent spatio-temporal structure from sparse sensor observations by aligning inference with local domains of dependence. This locality-aware design enables accurate reconstruction while avoiding the oversmoothing of global models and the scalability limitations of dense discretizations.

% An important practical strength of FieldFormer is parameter efficiency. In the shallow-water equation setting, FieldFormer uses approximately 71K trainable parameters, compared to 264K for SIREN, 284K for Fourier-MLP, and approximately 8M for SVGP. FieldFormer is thus the smallest among all neural field baselines, more than four times smaller than SIREN and Fourier-MLP and orders of magnitude smaller than SVGP. This demonstrates that the observed performance gains are not driven by increased model capacity, but by architectural inductive biases—specifically, adaptive neighborhood selection and local attention.

% \noindent\textbf{Limitations: Global Geometry and Heterogeneous Dynamics.}
% A key limitation of the current architecture is the use of globally constant velocity-scaling parameters, which impose a single anisotropic geometry across the entire domain. While this is effective when dominant directions of information flow are relatively uniform, it cannot fully capture spatially or temporally varying dynamics induced by heterogeneous forcing or nonstationary processes.

\section{Discussion and Conclusion}
\label{sec:discussion_conclusion}

This work studies sparse spatio-temporal reconstruction under the limits imposed by persistent sensor networks. Rather than treating sparse observations as sufficient for exact global field recovery, we argue that extreme spatial sparsity induces an observability gap: many globally distinct latent fields may remain consistent with the same sensor-space measurements. In this setting, reconstruction quality depends not only on model capacity, but also on the alignment between the model's inductive bias and the information geometry induced by the sensor network.

FieldFormer is designed around this perspective. By operating directly on coordinate-value observations, it avoids dependence on fixed grids or static graphs and supports mesh-free inference over irregular sensing layouts. Its adaptive local neighborhoods and local transformer aggregation align reconstruction with localized domains of dependence, making the model particularly effective for sensor-space imputation over persistent sparse sensor networks. Across synthetic PDE benchmarks and real-world monitoring datasets, FieldFormer performs strongly in this regime, especially when the target is to reconstruct missing or withheld measurements over an existing sensing topology.

Our results also clarify the distinction between sensor-space imputation and global field reconstruction. Sensor-space imputation is comparatively well grounded because predictions are made on the observational support of the deployed sensors. In contrast, global reconstruction away from observed support is substantially more prior-dependent. Across datasets, different methods perform best in different global reconstruction or sensor-holdout settings, with no universally dominant inductive bias. This variability supports the central claim of the paper: under extreme sparse sensing, global field reconstruction should be interpreted as prior-guided surrogate completion rather than identifiable recovery of the true latent physical state.

The ablation studies further highlight the role of FieldFormer's architectural choices. Replacing the local transformer with an MLP preserves the benefit of adaptive neighborhood construction but is less reliable on coupled or transport-dominated systems, suggesting that relational aggregation within local neighborhoods is important. Adding PINN-style physics regularization does not consistently improve performance and can degrade sensor-space accuracy, indicating that physics losses may conflict with sparse noisy observations when the residual objective is weakly constrained. Coordinate-dependent Gamma Fields improve selected metrics in heterogeneous regimes, but their gains are not uniform, suggesting that local geometry adaptation is a useful extension rather than a universally better default.

FieldFormer also has practical advantages for real-world sensing systems. Its mesh-free formulation supports irregular sensors, multi-resolution queries, and evolving sensor networks without requiring grid redesign or graph reconstruction. This is important in environmental monitoring, atmospheric sensing, pollution tracking, traffic infrastructure, industrial IoT, and scientific sensor deployments, where sensors may fail, move, or be added over time. In such settings, the ability to update the observational support while retaining the same coordinate-based model structure is a useful property for continual monitoring.

At the same time, several limitations remain. First, FieldFormer relies on local observability: when relevant local domains of dependence are not covered by the sensor network, reconstruction becomes underconstrained. Second, physics-based regularization remains delicate under sparse observations, especially when forcing terms, boundary conditions, or physical parameters are uncertain. Finally, FieldFormer does not eliminate the fundamental non-identifiability of global reconstruction under extreme sparse sensing; it provides a structured surrogate model whose reliability depends on sensor coverage and inductive-bias alignment.

These limitations suggest several directions for future work. Uncertainty-aware reconstruction could help distinguish well-supported sensor-space predictions from weakly constrained global completions. Adaptive sensor placement could improve observability by targeting regions where local dependencies are currently missed. Hybrid local-global architectures may combine FieldFormer's locality-aware inference with global latent priors, improving performance in regimes where local support is incomplete. Finally, more robust physics integration is needed to use partial governing knowledge without forcing models toward biased or overconstrained solutions.

Overall, FieldFormer provides both a practical method for sparse sensor-space imputation and a conceptual perspective on spatio-temporal reconstruction under sparse observability. The key lesson is that physical field reconstruction from sparse sensors is not merely an architectural problem; it is an information-constrained inference problem. By aligning model structure with the localized observational support of real sensor networks, FieldFormer offers a scalable and effective approach for constructing useful surrogate fields under severe sensing limitations.

\section*{Broader Impact}

Sparse sensor networks are widely used in environmental monitoring, atmospheric sensing, pollution tracking, traffic systems, industrial IoT, and scientific field studies, where dense measurement is often costly or infeasible. FieldFormer can improve the utility of such networks by imputing missing measurements and constructing useful surrogate fields from sparse, noisy observations. At the same time, our results emphasize that global reconstruction under extreme sparsity is not identifiable in general; model outputs should therefore be interpreted as prior-guided estimates rather than ground truth away from sensor support. Responsible deployment requires validation against held-out sensors where possible, uncertainty-aware reporting, and clear communication of where predictions are observationally supported versus weakly constrained.

\section*{Acknowledgements}
Ankit Bhardwaj and Lakshminarayanan Subramanian were supported by the NSF Grant (Award Number 2335773) titled "EAGER: Scalable Climate Modeling using Message-Passing Recurrent Neural Networks". Lakshminarayanan Subramanian was also funded in part by the NSF Grant (award number OAC-2004572) titled  "A Data-informed Framework for the Representation of Sub-grid Scale Gravity Waves to Improve Climate Prediction".